\documentclass[11pt]{article}
\usepackage[a4paper,margin=25mm]{geometry}
\usepackage{amsmath,amssymb,amsthm}
\usepackage[hidelinks]{hyperref}
\emergencystretch=2em
\newtheorem{theorem}{Theorem}
\title{A graphon does not determine a first-order limit\\
\large Disproof of Conjecture 00000009183}
\author{GitHub: gaochengzhecpu}\date{}
\begin{document}\maketitle
\begin{abstract}
The complete graphs and the complete graphs with one edge deleted have
the same constant-one graphon limit. Both sequences have first-order
limits, but the sentence asserting completeness distinguishes them at
every size. Thus graphons cannot characterize first-order limits of dense
graphs in the sense of determining their limiting first-order theories.
The ordinary dense graph limit theorem is unaffected.
\end{abstract}

\section*{Statement and interpretation}
The English source begins ``Definition: First-order limits of dense
graphs'' and describes graphons as the characterization of the limit
objects. The Chinese source likewise explicitly says first-order limits.
We test the following necessary consequence of a graphon characterization:
if two first-order convergent dense graph sequences have the same graphon
limit, they have the same eventual truth value for every first-order
sentence in the language with adjacency and equality.

The source does not specify a formal convergence definition. We use the
usual first-order convergence in which the proportions of satisfying
assignments converge for every formula; for a sentence these proportions
are its truth values, zero or one. Our sequences are convergent even in
this stronger formula sense, as shown below. If ``first-order'' in the
source were removed and only convergence of homomorphism densities were
intended, the familiar graphon characterization would be correct. The
counterexample concerns the first-order assertion actually present in both
versions, not that corrected assertion.

\section*{The counterexample}
For every $N\geq 2$, let $A_N=K_N$, and obtain $B_N$ by deleting exactly
one edge, whose endpoints are denoted $u,v$. Both graphs have $N$ vertices.
Their ordered edge densities are
\[
 \frac{N(N-1)}{N^2}\longrightarrow1,
 \qquad \frac{N(N-1)-2}{N^2}\longrightarrow1.
\]
In particular both sequences are dense.

Partition $[0,1)$ into $N$ intervals $I_i$ of length $1/N$, and let
$W_G$ be the usual adjacency step graphon: on $I_i\times I_j$ it is
$1$ if $i,j$ are adjacent and $0$ otherwise. Values at the endpoint
$1$ are immaterial. There are $N$ diagonal zero squares for $A_N$,
and $N+2$ zero squares for $B_N$. Consequently
\[
 \|1-W_{A_N}\|_1=\frac1N,
 \qquad \|1-W_{B_N}\|_1=\frac{N+2}{N^2}.
\]
The cut norm is bounded above by the $L^1$ norm, so both graphons
converge in cut norm, and hence in cut distance, to the same graphon
$W\equiv1$. This includes the diagonal error; it is not merely a
comparison of the two edge counts.

For completeness the same conclusion can be checked against every finite
test graph. If $F$ has $m$ edges, independent uniform choices of its
vertex images give, by the union bound,
\[
 0\leq 1-t(F,A_N)\leq\frac mN,
 \qquad 0\leq1-t(F,B_N)\leq\frac{m(N+2)}{N^2}.
\]
Each test edge fails with the displayed missing ordered-pair probability;
independence between the different failure events is not needed.
Thus all finite homomorphism densities converge to those of $W=1$.

Now consider the first-order sentence
\[
 \sigma:=\forall x\,\forall y\,(x\ne y\ \Longrightarrow\ E(x,y)).
\]
Every $A_N$ satisfies $\sigma$, whereas every $B_N$ fails it, witnessed
by $u,v$. Their limiting first-order theories therefore differ.

\section*{Existence of the first-order limits}
Here we also verify that the example is not relying on a sequence which
fails to have a first-order limit. For $A_N,A_M$, the duplicator wins the
$q$-round Ehrenfeucht--Fraisse game whenever $N,M\geq q$: equality of
already selected vertices is preserved, and a fresh vertex is matched by
a fresh vertex. All distinct selected vertices are adjacent.

For $B_N,B_M$, the duplicator preserves the pair of exceptional vertices
$\{u,v\}$, matching exceptional vertices to exceptional vertices and
ordinary vertices to ordinary vertices. Once an exceptional vertex is
selected its match is fixed, and the other exceptional vertex is matched
to the remaining exceptional vertex if selected. Fresh ordinary vertices
are always available when $N-2,M-2\geq q$. This is a partial isomorphism
after every move: the only missing edge between distinct vertices is the
exceptional pair. Thus every sentence of quantifier rank at most $q$ has
constant truth value on each sequence once $N\geq q+2$.

For the version of first-order convergence including formulas with free
variables, fix a formula with $k$ free variables and quantifier rank $q$.
The proportion of assignments with a repeated vertex or an exceptional
vertex is at most
\[
 \binom{k}{2}/N+2k/N,
\]
which tends to zero. (For $A_N$, omit the second term.) All remaining
assignments are injective and ordinary. Any two such assignments, in
graphs of sizes at least $k+q+2$, extend the same partial isomorphism
in the $q$-round game above. The formula therefore has one fixed truth
value on these assignments, independent of the sufficiently large size.
Its satisfaction proportion converges to that value. This proves the
claimed first-order convergence for every formula.

\section*{Formal verification and scope}
The Lean development uses Mathlib's \texttt{SimpleGraph} on
$\operatorname{Fin}(n+2)$, so $N=n+2$. It defines both graph sequences,
their actual zero--one adjacency kernels, and the explicit semantic
property represented by $\sigma$. It proves the truth and falsity of
that property on the two graphs for every $n$.

To connect the finite kernels to cut convergence, let $a_i,b_j\in[0,1]$
be arbitrary. Lean proves nonnegativity and the uniform bound
\[
 0\leq \frac1{N^2}\sum_{i,j}a_i b_j(1-W_{B_N}(i,j))
       \leq \frac{N+2}{N^2},
\]
and the same upper bound for $A_N$, together with convergence of the
upper bound to zero. This quantifies over all fractional rectangle
weights, not only subsets of the vertex set. For measurable $S,T$, put
$a_i=N\lambda(S\cap I_i)$ and $b_j=N\lambda(T\cap I_j)$.
The left side is exactly
$\int_{S\times T}(1-W_G)$ by the elementary integral formula for step
functions. Since $1-W_G\geq0$, these inequalities are precisely uniform
control of all cut rectangles.

The development defines this finite-kernel cut convergence and the
necessary assertion that a common constant-one limit determines the
eventual truth of $\sigma$, and proves its negation. The final theorem is
\begin{center}\texttt{Conjecture9183.conjecture9183\_false}.\end{center}
Its statement is the negation of
\begin{center}\texttt{GraphonDeterminesCompleteness}.\end{center}
The only axioms reported are \texttt{propext}, \texttt{Classical.choice},
and \texttt{Quot.sound}.

The identification of the finite weighted sums with measurable
step-graphon integrals, the homomorphism-density union bound, and the
Ehrenfeucht--Fraisse arguments are proved here on paper. Lean does not
implement a general syntax or satisfaction relation for arbitrary
first-order formulas, or a general graphon measure space. Its sentence
predicate is exactly the semantics of the displayed sentence. The source
is informal; the Lean proposition explicitly records the necessary
consequence being refuted instead of claiming a literal formalization of
undefined prose.

\section*{Source and reproduction}
The unmodified bilingual statement is supplied as \texttt{SOURCE.md}.
The source was frozen by the coordinating workflow; its SHA-256 is in the
provenance file. Lean is pinned to 4.19.0 and Mathlib commit
\begin{center}\texttt{c44e0c8ee63ca166450922a373c7409c5d26b00b}.\end{center}
From the supplied Lean project, run
\begin{center}\texttt{lake env lean -DwarningAsError=true Main.lean}.\end{center}
This draft was developed by a delegated AI agent with a self-review;
the coordinating agent reviews it separately. No external independent
review is claimed.
\end{document}
